\documentclass[11pt]{article}
\usepackage{preambulo}
\addbibresource{bibliografia.bib}

\begin{document}

\portada{Preprocesamiento de la instancia}

\tableofcontents
\newpage

%============================================================================
\section{Introducción y objetivo del preprocesamiento}
%============================================================================

El presente documento describe el preprocesamiento aplicado al conjunto de datos SUPPORT2
\parencite{support2}, correspondiente al estudio clínico sobre pronóstico de pacientes
hospitalizados con enfermedades graves descrito por \textcite{knaus1995support}. El objetivo del
procedimiento consistió en transformar el archivo original, que presenta valores ausentes,
variables categóricas en texto y columnas no utilizables, en dos conjuntos numéricos y completos
aptos para el entrenamiento y la evaluación de un modelo de clasificación binaria.

La variable objetivo seleccionada fue \texttt{hospdead}, que indica si el paciente falleció
durante la hospitalización analizada. La elección de esta variable condiciona la totalidad del
preprocesamiento, dado que determina cuáles columnas constituyen información legítima de entrada
y cuáles anticipan el desenlace que se pretende predecir.

El procedimiento se implementó en once etapas secuenciales e independientes, denominadas lotes,
cada una de las cuales recibe el resultado de la anterior y produce un archivo intermedio. Esta
organización permite ejecutar o revisar cualquier etapa de manera aislada sin necesidad de
recalcular las precedentes. La Figura \ref{fig:flujo} presenta la estructura general del
procedimiento.

\begin{figure}[H]
\centering
\begin{tikzpicture}[
  node distance=0.42cm,
  etapa/.style={rectangle, rounded corners=2pt, draw=gray!60, fill=gray!8,
                text width=8.2cm, minimum height=0.72cm, align=left, font=\small},
  grupo/.style={rectangle, rounded corners=3pt, draw=acento!70, thick, inner sep=4pt},
  flecha/.style={-{Latex[length=2mm]}, gray!70}
]
\node[etapa] (l1) {\textbf{1.} Descarga del conjunto de datos};
\node[etapa, below=of l1] (l2) {\textbf{2.} Auditoría inicial y estandarización de nombres};
\node[etapa, below=of l2, fill=acento!10] (l3) {\textbf{3.} Selección de variables: fuga y redundancia};
\node[etapa, below=of l3] (l4) {\textbf{4.} Eliminación de registros duplicados};
\node[etapa, below=of l4, fill=acento!10] (l5) {\textbf{5.} Partición estratificada 80/20};
\node[etapa, below=of l5] (l6) {\textbf{6.} Columnas con exceso de valores ausentes};
\node[etapa, below=of l6] (l7) {\textbf{7.} Acotamiento de valores atípicos};
\node[etapa, below=of l7] (l8) {\textbf{8.} Imputación de valores faltantes};
\node[etapa, below=of l8] (l9) {\textbf{9.} Codificación de variables categóricas};
\node[etapa, below=of l9] (l10) {\textbf{10.} Ordenamiento y escritura de los archivos};
\node[etapa, below=of l10] (l11) {\textbf{11.} Validación del resultado};

\foreach \a/\b in {l1/l2, l2/l3, l3/l4, l4/l5, l5/l6, l6/l7, l7/l8, l8/l9, l9/l10, l10/l11}
  \draw[flecha] (\a) -- (\b);

\node[anchor=west, font=\scriptsize\itshape, text=acento, align=left]
  at ([xshift=3mm]l3.east) {Determina qué\\información es\\admisible};
\node[anchor=west, font=\scriptsize\itshape, text=acento, align=left]
  at ([xshift=3mm]l5.east) {A partir de aquí\\todo se calcula\\sólo con el 80\%};
\end{tikzpicture}
\caption{Estructura secuencial del preprocesamiento. Las etapas destacadas corresponden a las dos
decisiones que condicionan la validez del resultado.}
\label{fig:flujo}
\end{figure}

La organización del presente documento sigue las tres dimensiones solicitadas: limpieza
(Sección \ref{sec:limpieza}), transformación (Sección \ref{sec:transformacion}) y reducción de la
dimensionalidad (Sección \ref{sec:reduccion}). No obstante, dado que la selección de variables
por fuga de información constituye la decisión de mayor incidencia sobre el resultado final y
antecede cronológicamente a las restantes etapas, se presenta de manera independiente en la
Sección \ref{sec:fuga}.

%============================================================================
\section{Descripción del conjunto de datos}
%============================================================================

El conjunto de datos SUPPORT2 proviene del estudio \textit{Study to Understand Prognoses and
Preferences for Outcomes and Risks of Treatments}, desarrollado en cinco hospitales de los Estados
Unidos con el propósito de caracterizar el pronóstico de pacientes hospitalizados con
enfermedades graves. El archivo original está conformado por 9{,}105 registros y 45 columnas,
donde cada registro corresponde a un paciente hospitalizado.

Los pacientes se distribuyen en ocho grupos diagnósticos: insuficiencia respiratoria aguda con
sepsis, insuficiencia cardíaca congestiva, enfermedad pulmonar obstructiva crónica, cirrosis,
cáncer de colon, cáncer de pulmón, coma y fallo multiorgánico con neoplasia.

Las variables disponibles pueden agruparse en cinco categorías: datos demográficos y de contexto
social, información diagnóstica, puntajes de severidad clínica medidos al tercer día del estudio,
signos vitales y resultados de laboratorio, e indicadores de funcionalidad en actividades de la
vida diaria. La Tabla \ref{tab:variables} presenta la totalidad de las variables con el rol que les
asigna la ficha oficial del repositorio, su descripción, tipo y proporción de valores ausentes.

\begin{longtable}{p{1.7cm} p{5.1cm} p{1.4cm} p{1.6cm} r}
\caption{Variables del conjunto de datos original. Descripciones tomadas de la ficha oficial del repositorio UCI.}
\label{tab:variables}\\
\toprule
\textbf{Variable} & \textbf{Descripción} & \textbf{Rol} & \textbf{Tipo} & \textbf{\% nulos} \\
\midrule
\endfirsthead
\toprule
\textbf{Variable} & \textbf{Descripción} & \textbf{Rol} & \textbf{Tipo} & \textbf{\% nulos} \\
\midrule
\endhead
\bottomrule
\endfoot
age & Edad del paciente en años & Entrada & Numérica & 0.0 \\
sex & Sexo del paciente & Entrada & Binaria & 0.0 \\
dzgroup & Subcategoría de enfermedad del paciente & Entrada & Numérica & 0.0 \\
dzclass & Categoría de enfermedad del paciente & Entrada & Numérica & 0.0 \\
num\_co & Número de enfermedades simultáneas o comorbilidades & Entrada & Numérica & 0.0 \\
edu & Años de educación & Entrada & Numérica & 17.9 \\
income & Ingreso del paciente & Entrada & Numérica & 32.8 \\
scoma & Puntaje de coma SUPPORT al día 3, basado en la escala de Glasgow (estimado por un modelo) & Entrada & Numérica & 0.0 \\
charges & Cobro hospitalario & Entrada & Numérica & 1.9 \\
totcst & Costo total según la razón costo-cobro (RCC) & Entrada & Numérica & 9.8 \\
totmcst & Costo total por micro-costeo & Entrada & Numérica & 38.2 \\
avtisst & Puntaje TISS promedio de los días 3 al 25, método de costeo de intervenciones terapéuticas & Entrada & Numérica & 0.9 \\
race & Raza del paciente & Entrada & Numérica & 0.5 \\
sps & Puntaje de fisiología SUPPORT al día 3 (estimado por un modelo) & Entrada & Numérica & 0.0 \\
aps & Puntaje de fisiología APACHE III al día 3 & Entrada & Numérica & 0.0 \\
surv2m & Estimación de supervivencia a 2 meses del modelo SUPPORT al día 3 & Entrada & Numérica & 0.0 \\
surv6m & Estimación de supervivencia a 6 meses del modelo SUPPORT al día 3 & Entrada & Numérica & 0.0 \\
hday & Día de hospitalización en que el paciente ingresó al estudio & Entrada & Numérica & 0.0 \\
diabetes & Presencia de diabetes como comorbilidad & Entrada & Binaria & 0.0 \\
dementia & Presencia de demencia como comorbilidad & Entrada & Binaria & 0.0 \\
ca & Presencia de cáncer: sí, metastásico o no & Entrada & Numérica & 0.0 \\
prg2m & Estimación del médico sobre la supervivencia a 2 meses & Entrada & Numérica & 18.1 \\
prg6m & Estimación del médico sobre la supervivencia a 6 meses & Entrada & Numérica & 17.9 \\
dnr & Existencia de orden de no reanimar & Entrada & Numérica & 0.3 \\
dnrday & Día de la orden de no reanimar; negativo si es previa al estudio & Entrada & Numérica & 0.3 \\
meanbp & Presión arterial media, medida al día 3 & Entrada & Numérica & 0.0 \\
wblc & Recuento de glóbulos blancos en miles, medido al día 3 & Entrada & Numérica & 2.3 \\
hrt & Frecuencia cardíaca, medida al día 3 & Entrada & Numérica & 0.0 \\
resp & Frecuencia respiratoria, medida al día 3 & Entrada & Numérica & 0.0 \\
temp & Temperatura en grados Celsius, medida al día 3 & Entrada & Numérica & 0.0 \\
pafi & Razón PaO$_2$/FiO$_2$, medida al día 3 & Entrada & Numérica & 25.5 \\
alb & Albúmina sérica, medida al día 3 & Entrada & Numérica & 37.0 \\
bili & Bilirrubina, medida al día 3 & Entrada & Numérica & 28.6 \\
crea & Creatinina sérica, medida al día 3 & Entrada & Numérica & 0.7 \\
sod & Concentración sérica de sodio, medida al día 3 & Entrada & Numérica & 0.0 \\
ph & pH de la sangre arterial & Entrada & Numérica & 25.1 \\
glucose & Glucosa, medida al día 3 & Entrada & Numérica & 49.4 \\
bun & Nitrógeno ureico en sangre, medido al día 3 & Entrada & Numérica & 47.8 \\
urine & Volumen de orina, medido al día 3 & Entrada & Numérica & 53.4 \\
adlp & Índice de actividades de la vida diaria, respondido por el paciente & Entrada & Numérica & 62.0 \\
adls & Índice de actividades de la vida diaria, respondido por un allegado & Entrada & Numérica & 31.5 \\
adlsc & Índice de actividades de la vida diaria imputado y calibrado al allegado & Entrada & Numérica & 0.0 \\
death & Fallecimiento en cualquier momento hasta el registro del National Death Index al 31 de diciembre de 1994 & Objetivo & Binaria & 0.0 \\
hospdead & Fallecimiento durante la hospitalización & Objetivo & Binaria & 0.0 \\
sfdm2 & Nivel de discapacidad funcional en escala de 1 a 5 & Objetivo & Numérica & 15.4 \\
\end{longtable}


El conjunto de datos contiene tres variables susceptibles de ser utilizadas como objetivo:
\texttt{death}, que registra el fallecimiento en cualquier momento del seguimiento;
\texttt{hospdead}, que registra el fallecimiento durante la hospitalización analizada; y
\texttt{sfdm2}, que corresponde al nivel de discapacidad funcional evaluado en una entrevista
posterior al alta. Se seleccionó \texttt{hospdead} por tratarse de un desenlace delimitado
temporalmente y observable para la totalidad de los registros.

La Figura \ref{fig:clases} muestra la distribución de la variable objetivo en el conjunto de
entrenamiento. Se observa un desbalance de clases en el que los fallecimientos constituyen
aproximadamente una cuarta parte de los casos, condición que deberá considerarse durante la
evaluación del modelo.

\begin{figure}[H]
\centering
\includegraphics[width=0.72\textwidth]{clases.png}
\caption{Distribución de la variable objetivo \texttt{hospdead} en el conjunto de entrenamiento.}
\label{fig:clases}
\end{figure}

%============================================================================
\section{Auditoría inicial}
%============================================================================

Previo a cualquier modificación del conjunto de datos se realizó una auditoría destinada a
establecer el estado inicial de la información. Este diagnóstico constituye el punto de
referencia contra el cual se verifican posteriormente los efectos de cada etapa.

La lectura del archivo requirió tres operaciones de normalización previas. En primer lugar, se
declararon explícitamente los marcadores que deben interpretarse como valores ausentes
(\texttt{?}, \texttt{NA}, \texttt{N/A}, cadena vacía y \texttt{null}), dado que de lo contrario
serían tratados como texto válido y distorsionarían el conteo de faltantes. En segundo lugar, se
estandarizaron los nombres de las columnas a minúsculas con guion bajo, debido a que
identificadores como \texttt{num.co} contienen un punto que interfiere con la notación de acceso
a atributos. Finalmente, se descartaron los registros completamente vacíos.

A partir del conjunto normalizado se construyó una tabla de diagnóstico que reporta, para cada
columna, su tipo de dato, el número y porcentaje de valores ausentes, y la cantidad de valores
distintos. El cálculo del porcentaje de faltantes se apoya en que la media de una serie booleana
equivale a la proporción de casos verdaderos, por lo que el promedio de los indicadores de
ausencia expresa directamente dicha proporción.

La auditoría reveló una distribución marcadamente heterogénea de valores ausentes. Mientras que
variables como la edad, el sexo o el número de comorbilidades se encuentran completas, otras
superan el cincuenta por ciento de datos faltantes, situación que condiciona su posterior
tratamiento y motiva las decisiones documentadas en la Sección \ref{sec:limpieza}.

%============================================================================
\section{Selección de variables y fuga de información}
\label{sec:fuga}
%============================================================================

\subsection{Definición del problema}

Se denomina fuga de información a la presencia, dentro de las variables de entrada, de datos que
no estarían disponibles en el momento en que el modelo debe emitir su predicción. Un modelo
entrenado bajo esta condición alcanza un desempeño elevado durante la evaluación, pero dicho
desempeño no es reproducible en condiciones reales de uso, dado que la información que lo
sustenta sólo existe una vez conocido el desenlace.

En el caso analizado, la fuga adopta dos formas distintas que se abordan por separado. La primera
corresponde a columnas que contienen, de manera directa o indirecta, la respuesta que se pretende
predecir; su tratamiento se documenta en la presente sección. La segunda deriva del orden en que
se ejecutan las operaciones del preprocesamiento y se documenta en la Sección \ref{sec:orden}.

La identificación de las columnas afectadas se realizó mediante el análisis de la ficha oficial
del conjunto de datos, la revisión del momento en que cada variable se registra durante la
hospitalización, y la verificación empírica de su relación con la variable objetivo.

\subsection{Variables que constituyen el desenlace}

Un primer conjunto de columnas corresponde a mediciones del mismo fenómeno que se desea predecir.

La variable \texttt{death} registra el fallecimiento del paciente en cualquier momento del
seguimiento, período que puede extenderse durante años. Dado que todo paciente que fallece durante
la hospitalización fallece necesariamente dentro del seguimiento, ambas variables responden
esencialmente a la misma pregunta.

La variable \texttt{sfdm2} consigna el nivel de discapacidad funcional evaluado mediante una
entrevista realizada aproximadamente dos meses después del ingreso. Dicha entrevista sólo pudo
realizarse a los pacientes que sobrevivieron hasta esa fecha, por lo que la mera presencia de un
valor en esta columna constituye un indicador de supervivencia.

Las variables \texttt{surv2m} y \texttt{surv6m} corresponden a la probabilidad de supervivencia
estimada por el modelo estadístico del propio estudio SUPPORT, mientras que \texttt{prg2m} y
\texttt{prg6m} recogen la estimación subjetiva formulada por el médico tratante. Su incorporación
como variables de entrada implicaría reproducir el resultado de un modelo previo o el juicio de un
especialista, en lugar de aprender a partir de las mediciones clínicas.

\subsection{Variables conocidas únicamente al alta}

Un segundo conjunto de columnas registra información generada durante la hospitalización, cuyo
valor definitivo sólo se conoce una vez concluido el episodio.

Las variables \texttt{charges}, \texttt{totcst} y \texttt{totmcst} corresponden al cobro y al
costo total de la hospitalización. Estas cifras se consolidan al momento del alta, de modo que no
se encuentran disponibles mientras el paciente permanece internado. La diferencia observada entre
grupos resulta considerable: el cobro promedio asciende a 92{,}638 unidades monetarias entre los
pacientes fallecidos frente a 48{,}683 entre los sobrevivientes, diferencia atribuible a la
duración y la intensidad del episodio.

La variable \texttt{avtisst} corresponde al promedio del puntaje de intervenciones terapéuticas
entre los días 3 y 25 del estudio. Al promediarse a lo largo de la estancia, su valor queda
determinado por la duración de esta última.

Las variables \texttt{dnr} y \texttt{dnrday} registran la orden de no reanimar y la fecha en que
fue suscrita. Esta orden no constituye un dato clínico de ingreso sino una decisión adoptada
durante la hospitalización, con frecuencia cuando el equipo médico ya ha constatado el deterioro
del paciente. La Tabla \ref{tab:dnr} evidencia esta relación.

\begin{table}[H]
\centering
\caption{Distribución de la orden de no reanimar según el desenlace hospitalario}
\label{tab:dnr}
\small
\begin{tabular}{p{4.5cm} r r}
\toprule
\textbf{Estado de la orden} & \textbf{Sobrevive} & \textbf{Fallece} \\
\midrule
Suscrita antes del ingreso  & 2.4\%  & 3.2\%  \\
Suscrita después del ingreso & 16.8\% & 77.9\% \\
Sin orden de no reanimar     & 80.8\% & 18.9\% \\
\bottomrule
\end{tabular}
\end{table}

Entre los pacientes que fallecen durante la hospitalización, el 77.9\% cuenta con una orden
suscrita con posterioridad al ingreso, frente al 16.8\% de quienes sobreviven. Esta asociación no
refleja una capacidad predictiva de la variable sino la anticipación del desenlace por parte del
personal médico.

Se conservó en cambio la variable \texttt{hday}, correspondiente al día de hospitalización en que
el paciente ingresó al estudio. Aunque presenta diferencias entre grupos, su valor se conoce en el
momento del ingreso y constituye por tanto información legítima.

\subsection{Variables redundantes}

Un tercer conjunto de columnas no introduce fuga de información, pero duplica contenido ya
presente en otras variables.

La tabulación cruzada entre \texttt{dzgroup} y \texttt{dzclass} evidencia que cada categoría de la
primera corresponde invariablemente a una única categoría de la segunda, sin excepciones. En
consecuencia, \texttt{dzclass} constituye un agrupamiento derivado que incrementa la
dimensionalidad sin aportar información adicional. Se conservó \texttt{dzgroup} por presentar
mayor detalle.

De manera análoga, la variable \texttt{adlsc} corresponde, conforme a la documentación oficial, a
una versión calibrada e imputada de \texttt{adls}. Mantener ambas implicaría duplicar la misma
información y, adicionalmente, imputar el 31\% de valores ausentes que presenta la versión
original. Se conservó \texttt{adlsc}.

Se evaluó asimismo la relación entre \texttt{sps} y \texttt{aps}, ambos puntajes de severidad
medidos al tercer día. La correlación observada resultó elevada pero inferior a la unidad, por lo
que no constituyen mediciones equivalentes y ambas fueron conservadas.

\subsection{Resumen de las exclusiones}

La Tabla \ref{tab:exclusiones} consigna la totalidad de las columnas excluidas junto con el motivo
correspondiente.

\begin{table}[H]
\centering
\caption{Columnas excluidas del conjunto de modelado}
\label{tab:exclusiones}
\small
\begin{tabular}{p{2.2cm} p{8.5cm}}
\toprule
\textbf{Columna} & \textbf{Motivo de exclusión} \\
\midrule
\texttt{death} & otra variable de muerte, casi la misma pregunta que hospdead \\
\texttt{sfdm2} & otra variable objetivo, medida despues del alta \\
\texttt{surv2m} & estima el desenlace: supervivencia a 2 meses segun el modelo SUPPORT \\
\texttt{surv6m} & estima el desenlace: supervivencia a 6 meses segun el modelo SUPPORT \\
\texttt{prg2m} & estima el desenlace: supervivencia a 2 meses segun el medico tratante \\
\texttt{prg6m} & estima el desenlace: supervivencia a 6 meses segun el medico tratante \\
\texttt{dzclass} & redundante: es un resumen derivado de dzgroup \\
\texttt{adls} & redundante: adlsc ya contiene la misma informacion, sin huecos \\
\texttt{dnr} & la orden de no reanimar suele firmarse cuando el paciente ya esta empeorando \\
\texttt{dnrday} & dia de la orden de no reanimar: misma razon que dnr \\
\texttt{charges} & cobro total de la hospitalizacion, se conoce recien al dar de alta \\
\texttt{totcst} & costo total de la hospitalizacion, se conoce recien al dar de alta \\
\texttt{totmcst} & costo total detallado, se conoce recien al dar de alta \\
\texttt{avtisst} & puntaje TISS promedio dias 3-25, que la ficha define como metodo de costeo \\
\bottomrule
\end{tabular}
\end{table}


\subsection{Verificación empírica del efecto}
\label{sec:evidencia}

Con el propósito de cuantificar la incidencia de las exclusiones, se entrenó una regresión
logística bajo dos configuraciones idénticas en todo excepto en la presencia de las seis columnas
posteriores al ingreso. Los resultados se presentan en la Tabla \ref{tab:fuga}.

\begin{table}[H]
\centering
\caption{Efecto de las variables posteriores al ingreso sobre el desempeño del modelo}
\label{tab:fuga}
\small
\begin{tabular}{p{6.5cm} r r}
\toprule
\textbf{Configuración} & \textbf{Columnas} & \textbf{AUC} \\
\midrule
Con las variables posteriores al ingreso & 48 & 0.948 \\
Sin las variables posteriores al ingreso & 41 & 0.879 \\
\bottomrule
\end{tabular}
\end{table}

La diferencia observada resulta determinante para la interpretación del resultado. Un área bajo la
curva de 0.948 en la predicción de mortalidad hospitalaria excede el desempeño reportado para los
modelos publicados sobre este mismo conjunto de datos, los cuales se sitúan entre 0.80 y 0.86. El
valor de 0.879 obtenido tras la exclusión se encuentra en cambio dentro del rango esperable.

La reducción del desempeño no debe interpretarse por tanto como una pérdida de capacidad
predictiva, sino como la eliminación de una ventaja ilegítima: el modelo inicial no predecía el
fallecimiento a partir de información clínica, sino que lo inferían a partir de registros
generados como consecuencia del propio desenlace.


%============================================================================
\section{Limpieza}
\label{sec:limpieza}
%============================================================================

\textit{Sección en elaboración.}

%============================================================================
\section{Transformación}
\label{sec:transformacion}
%============================================================================

\textit{Sección en elaboración.}

%============================================================================
\section{Reducción de la dimensionalidad}
\label{sec:reduccion}
%============================================================================

\textit{Sección en elaboración.}

%============================================================================
\section{Fuga de información por orden de operaciones}
\label{sec:orden}
%============================================================================

\textit{Sección en elaboración.}

%============================================================================
\section{Validación del resultado}
\label{sec:validacion}
%============================================================================

\textit{Sección en elaboración.}

%============================================================================
\section{Conclusiones}
\label{sec:conclusiones}
%============================================================================

\textit{Sección en elaboración.}

\printbibliography[title={Referencias}]

\end{document}
